\section{Uniform regulator estimates}\label{app:regulator-estimates}

\subsection{Frequency blocks}

Work first on a finite Brown--Henneaux block $m,n\leq M$. After applying the PSL section, every profile vanishes at both interval anchors. The finite-action boundary contribution decomposes into

\begin{align}
\mathcal F_{M,L,\epsilon,R}
=\mathcal F_{\rm trans}
+\mathcal F_{\rm wall}
+\mathcal F_{\rm joint}
+\mathcal F_{\rm outer\,source}
+\mathcal F_{\rm outer\,corner}
+\mathcal F_{\rm outer\,joint}.
\end{align}

The individual estimates are

\[\begin{array}{c|c}
\text{sector} & \text{uniform block bound}\\ \hline
\text{transition} & C M^6L^{-2}\\
\text{complete compensated wall} & C M^7L^{-2}\\
\text{far wall tail} & C M^7(yL^{-5}+L^{-6})\\
\text{far Hayward joint} & C M^6yR^{-5}\\
\text{outer Brown--York source} & 0\\
\text{outer Cauchy corner} & C(M^6L^2R^{-2}+M^6L^{-2})\\
\text{outer Hayward mismatch} & C M^6L^2R^{-2}
\end{array} \label{eq:regulator-table}
\]

where $y=\tanh(\epsilon/2)$. The complete wall estimate already contains its endpoint shift; adding a separate anchor-shift term would double count the same variation.

\subsection{Explicit diagonal schedule}

For

\begin{align}
L=M^{10},
\qquad
y=\frac1{2M^{10}},
\qquad
R=M^{30},
\end{align}

the entries of (\ref{eq:regulator-table}) scale as

\[\begin{array}{c|c}
\text{sector} & \text{rate}\\ \hline
\text{transition} & O(M^{-14})\\
\text{complete compensated wall} & O(M^{-13})\\
\text{far wall tail} & O(M^{-53})\\
\text{far joint} & O(M^{-154})\\
\text{outer source} & 0\\
\text{outer corner} & O(M^{-34}+M^{-14})\\
\text{outer joint} & O(M^{-34}).
\end{array}
\]

The total flux is therefore $O(M^{-13})$. The precise exponents are not canonical; they provide one explicit admissible schedule. What matters is that the estimates are uniform in the block and that the regulator limits are taken together.

\subsection{Why the limits cannot be separated naively}

On the RT curve, a convenient proper profile is

\begin{align}
p_L(\phi)
=\frac{L^2\cos2\phi}{1+L^2\cos2\phi}.
\end{align}

In a frequency-$N$ anchor layer, $\phi=\pi/4-y/N$,

\begin{align}
p_L\left(\frac\pi4-\frac yN\right)
\longrightarrow
\frac{2y}{2y+N/L^2}.
\end{align}

Thus the local completion is recovered uniformly only if $M/L^2\to0$. Sending $M\to\infty$ at fixed $L$ switches off the completion in the high-frequency anchor layer, while a finite limit of $M/L^2$ leaves a schedule-dependent profile.

\subsection{Sobolev summation}

The separated endpoint chart contains the momentum trace

\begin{align}
p_\pm[f]
=\frac16\left[-f''(\pm a)-f(\pm a)\pm f'(\pm a)\right].
\end{align}

For a Fourier block its squared operator norm behaves as

\begin{align}
\sum_{m\leq M}m^{4-2\sigma}.
\end{align}

Uniform boundedness therefore holds exactly for $\sigma>5/2$. This is the sharp point-trace threshold for the separated chart. The cancellation in the combined form lowers the ordinary continuity threshold to $H^2$, as proved in Section \ref{sec:positivity}.

\subsection{Time transport}

Boundary time evolution multiplies each chiral frequency by a unit phase. Hence all polynomial frequency bounds in (\ref{eq:regulator-table}) are unchanged. Conjugating the endpoint and HW sections as in (\ref{eq:time-sections}) transports the whole finite-action stratum inventory. Since the total flux vanishes uniformly, the integrated symplectic forms on two transported slices agree.
